\chapter{Literature Review}
\label{ch:literature}

This chapter surveys the four bodies of work on which this thesis draws: the
clinical definition of kidney risk, machine-learning approaches to CKD
prediction, techniques for learning under class imbalance and for explaining
individual predictions, and the use of large language models---with retrieval
augmentation, parameter-efficient fine-tuning, and safety guardrails---in
patient-facing healthcare applications. It closes with an explicit statement of
the gap this research addresses.

\section{Clinical Background: Defining Kidney Risk}
\label{sec:clinical}

\subsection{The KDIGO risk grid}

The KDIGO 2024 clinical practice guideline \cite{kdigo2024} defines CKD as
abnormalities of kidney structure or function present for more than three
months. Risk is stratified on two orthogonal axes. The first is the estimated
glomerular filtration rate (\eGFR), computed from serum creatinine and
expressed in mL/min/1.73\,m$^2$, which is binned into categories G1 ($\geq$90)
through G5 ($<$15). The second is albuminuria, measured by the urine
albumin-to-creatinine ratio (\UACR) in mg/g, binned as A1 ($<$30), A2 (30--300)
and A3 ($>$300). The cross-product of these two axes yields the four risk
categories---low, moderately increased, high and very high---shown in
Figure~\ref{fig:kdigo} of Chapter~\ref{ch:methodology}.

This grid is the label source for the present work. It is important to be
precise about the direction of the dependency: \eGFR{} and \UACR{} define
\emph{what the model is trying to predict}, and are deliberately withheld from
\emph{what the model is allowed to see}.

\subsection{Modifiable risk factors}

The epidemiological literature and the CDC's public health guidance
\cite{cdcbasics,cdcprevent,cdcrisk} identify diabetes and hypertension as the
two dominant causes of CKD, together accounting for the majority of incident
cases, followed by obesity, cardiovascular disease, age, and family history.
The ADA/KDIGO consensus on diabetes management in CKD \cite{adakdigo2022} and
the KDIGO blood-pressure guideline \cite{kdigo2021bp} both emphasise that
glycaemic and blood-pressure control materially alter the trajectory of kidney
decline. The KDOQI nutrition guideline \cite{kdoqi2020} governs dietary advice
and is the source of an important constraint adopted in this work: individualised
numeric targets for protein, potassium, phosphorus or sodium must not be issued
to a patient without the involvement of a registered or renal dietitian.

The practical consequence for a screening model is favourable: the dominant
causal drivers of CKD---blood pressure, glycaemic markers, and adiposity---are
all routinely measured, and none of them require a kidney-specific laboratory
order.

\section{Machine Learning for CKD Prediction}
\label{sec:ml_ckd}

\subsection{The high-accuracy literature and its input dependency}

A substantial body of work applies supervised learning to CKD detection.
Chittora et al.\ \cite{chittora2021} compare a range of classifiers on the UCI
CKD dataset and report accuracies above 98\%. Aljaaf et al.\ \cite{aljaaf2019}
report comparable performance with predictive analytics support, and Yildirim
\cite{yildirim2017} examines multilayer perceptron performance on imbalanced
CKD data. These results are internally valid, but two properties limit their
applicability to community screening.

First, the feature sets almost universally contain serum creatinine, blood urea,
or \eGFR{} itself. Since the CKD label in these datasets is defined by kidney
function, a model given creatinine is effectively re-deriving a quantity that
the label already encodes. High reported accuracy is then expected, and it does
not demonstrate that the disease could have been detected \emph{before} the
confirmatory test was ordered.

Second, the widely used UCI CKD dataset contains only 400 records, with a
relatively balanced class distribution that does not reflect the roughly
10:1 prevalence imbalance found in an unselected adult population. Models
validated only on such data have not been tested against the imbalance that
dominates real screening.

\subsection{Prediction models for kidney-failure progression}

A separate and more clinically rigorous line of work predicts \emph{progression}
among patients already diagnosed with CKD. Ramspek et al.\ \cite{ramspek2021}
systematically review such models and find substantial heterogeneity in
predictor sets, frequent risk of bias, and limited external validation. The
canonical example, the Kidney Failure Risk Equation, requires \eGFR{} and
\UACR{} as inputs by design. These models answer a different question from the
one posed here: they estimate how fast a known patient will deteriorate,
whereas this thesis asks who should be tested in the first place.

\subsection{Positioning of the present work}

Table~\ref{tab:lit_ml} summarises the contrast between representative prior work
and the approach taken here.

\begin{table}[H]
\centering
\caption{Representative machine-learning work on CKD compared with this thesis.}
\label{tab:lit_ml}
\small
\begin{tabularx}{\textwidth}{@{}L{1.25in} L{0.85in} X L{0.95in}@{}}
\toprule
\textbf{Work} & \textbf{Data} & \textbf{Input features} & \textbf{Screening
usable pre-test?} \\
\midrule
Chittora et al.\ \cite{chittora2021} & UCI CKD ($n$=400) &
Includes serum creatinine, blood urea, specific gravity, albumin &
No --- requires the confirmatory assay \\
Aljaaf et al.\ \cite{aljaaf2019} & UCI CKD ($n$=400) &
Includes creatinine-derived markers & No \\
Yildirim \cite{yildirim2017} & UCI CKD ($n$=400) &
Includes laboratory renal markers & No \\
Ramspek et al.\ \cite{ramspek2021} (review) & Cohorts of diagnosed patients &
\eGFR{} and \UACR{} required by design & No --- assumes diagnosis \\
\textbf{This work} & \textbf{NHANES 2021--23 ($n$=5{,}552)} &
\textbf{Vitals, routine blood chemistry, comorbidity history; creatinine /
\eGFR{} / \UACR{} excluded} & \textbf{Yes} \\
\bottomrule
\end{tabularx}
\end{table}

\section{Learning under Class Imbalance}
\label{sec:imbalance}

In an unselected adult population the high-risk class is rare. A classifier
trained to minimise overall error can achieve high accuracy by predicting the
majority class almost everywhere, at the cost of failing precisely on the
minority class of interest. Lema\^{i}tre et al.\ \cite{lemaitre2017} provide the
standard toolbox for addressing this: random over-sampling of the minority
class, under-sampling of the majority, and hybrid schemes such as SMOTE-based
synthesis.

Three consequences guide the design in this thesis. First, aggregate accuracy is
an inadequate primary metric; balanced accuracy, per-class recall and the
confusion matrix are reported instead. Second, resampling should be applied
within the training pipeline only, so that evaluation reflects the true class
prior. Third---and central to the present design---not all misclassifications
are equally harmful: predicting ``low'' for a truly high-risk person removes
the very referral the system exists to generate, whereas predicting ``high''
for a low-risk person causes an unnecessary but ultimately corrective
confirmatory test. This asymmetry motivates decision thresholds tuned against a
specific error type rather than against symmetric loss.

\section{Cascaded and Hierarchical Classification}
\label{sec:cascade_lit}

Decomposing a multi-class problem into a sequence of binary decisions is a
long-standing technique in detection systems, where an early stage is tuned for
high sensitivity to pass a manageable candidate set to a more discriminative
later stage. The advantage relevant here is that each stage can be optimised,
and its operating threshold selected, against a different criterion: the first
stage for recall of the at-risk population, the second for separating severity
within that population. A single three-class model, by contrast, exposes only
one implicit decision boundary structure and offers no natural place to encode
the asymmetric cost described above.

\section{Explainability of Individual Predictions}
\label{sec:xai}

For a screening tool intended for non-expert users, an unexplained category is
of limited value. Two families of local explanation are widely used. SHAP
\cite{lundberg2017} computes Shapley-value attributions with strong theoretical
guarantees but requires either a model-specific efficient implementation or
substantial sampling; LIME \cite{ribeiro2016} fits a local surrogate model
around the instance. A third, simpler family is single-feature perturbation: the
value of one feature is replaced by a reference value and the change in
predicted probability is recorded.

The present work uses the perturbation approach, applied against the complete
deployed cascade rather than a surrogate, and reports it explicitly as such.
The trade-off is accepted deliberately: the method is fast enough to run inside
a request, requires no additional dependency, and is trivially auditable, but it
does not capture interactions between features and is not a Shapley
attribution. Chapter~\ref{ch:methodology} states the method formally and
Chapter~\ref{ch:discussion} returns to its limitations.

\section{Large Language Models in Healthcare}
\label{sec:llm_health}

Singhal et al.\ \cite{singhal2023} demonstrated that large language models
encode a substantial amount of clinical knowledge, achieving strong performance
on medical question-answering benchmarks, while also documenting failure modes
in which fluent output is factually wrong. Thirunavukarasu et al.\
\cite{thirunavukarasu2023} survey clinical applications and identify
hallucination, outdated knowledge, and the absence of provenance as the primary
barriers to deployment. The consistent finding across this literature is that
fluency and correctness are decoupled: an incorrect statement is delivered with
exactly the same confidence as a correct one, which is particularly hazardous
for a lay audience with no basis for scepticism. Broader reviews of natural
language processing over clinical text \cite{li2022nlp} reach a similar
conclusion regarding the need for explicit provenance.

For patient education specifically, the requirements are narrower than for
clinical decision support but no less strict. The output must not diagnose,
must not alter medication, must not issue individualised clinical targets, and
must escalate appropriately when presented with red-flag findings.

\section{Retrieval-Augmented Generation}
\label{sec:rag_lit}

Retrieval-augmented generation \cite{lewis2020} conditions a language model on
passages retrieved from an external corpus at inference time, rather than
relying solely on knowledge stored in the model's parameters. For a clinical
education application this offers three concrete advantages: the knowledge
source is explicit and citable, it can be updated by re-indexing rather than
retraining when guidelines change, and the retrieved passages constrain the
model's factual claims.

Retrieval quality depends on the embedding model. BGE-M3 \cite{chen2024bgem3}
provides multilingual, multi-granularity dense embeddings trained with
self-knowledge distillation, and is used in this work to embed both the
knowledge-base passages and the per-patient query.

A limitation of naive top-$k$ retrieval, observed directly during development
and addressed in Chapter~\ref{ch:methodology}, is that in a narrow domain the
same small set of generic passages tends to be returned for every query,
because all queries are semantically similar. The retrieval strategy adopted
here therefore incorporates the patient's own attributed risk factors into the
selection process rather than relying on similarity ranking alone.

\section{Parameter-Efficient Fine-Tuning}
\label{sec:peft_lit}

Full fine-tuning of a multi-billion-parameter model requires optimiser state and
gradients for every weight, which is beyond the memory budget of commonly
available hardware and risks degrading the model's general instruction-following
ability. Low-Rank Adaptation (LoRA) \cite{hu2022lora} freezes the pretrained
weights and injects a trainable low-rank decomposition into selected projection
matrices, so that the adapted weight becomes
\begin{equation}
W' = W + \frac{\alpha}{r} B A ,
\label{eq:lora_lit}
\end{equation}
where $W$ is the frozen weight, $A \in \mathbb{R}^{r \times d}$ and
$B \in \mathbb{R}^{d \times r}$ are the trainable factors, $r$ is the rank and
$\alpha$ a scaling constant. Only $A$ and $B$ receive gradients, reducing the
trainable parameter count by orders of magnitude.

QLoRA \cite{dettmers2023qlora} extends this by holding the frozen base weights
in 4-bit NormalFloat quantisation during training, further reducing memory
while keeping the adapter itself in higher precision. It is worth stating the
relationship precisely, because the terminology is frequently conflated: QLoRA
is not a different adapter type from LoRA. It is LoRA training performed with a
quantised frozen base. The artefact produced is an ordinary LoRA adapter, and it
may subsequently be attached to a base model held at any precision. Quantisation
is therefore a memory-saving choice rather than a methodological one: where the
GPU can hold the full-precision base, as in the present work, standard LoRA in
bfloat16 is preferable because it avoids adding quantisation error to the frozen
weights. This thesis uses standard LoRA, not QLoRA.

\section{Safety Guardrails for Generative Systems}
\label{sec:guardrails}

Three broad strategies exist for constraining generative output. \emph{Prompt
engineering} states the rules in the instruction; it is cheap and flexible but
provides no guarantee, since compliance is probabilistic. \emph{Fine-tuning}
attempts to internalise the rules in the weights; it improves the prior but
remains statistical and is difficult to verify exhaustively. \emph{Post-generation
validation} inspects the produced text and blocks, rewrites, or replaces it; it
is the only one of the three that can offer a deterministic guarantee for the
patterns it covers, at the cost of being unable to anticipate patterns it does
not.

An increasingly common fourth option is to use a second language model as a
judge. This scales to novel phrasings but reintroduces probabilistic behaviour
at exactly the point where a guarantee is wanted, and adds inference cost and
latency.

The literature reporting \emph{measured} unsafe-output rates for deployed
patient-facing systems---particularly rates attributed to individual
architectural components---is sparse. This observation directly motivates the
evaluation design in Chapter~\ref{ch:results}.

\section{Research Gap}
\label{sec:gap}

Synthesising the above, four gaps are identified.

\begin{enumerate}
  \item \textbf{Input circularity in CKD screening models.} Existing
        high-accuracy CKD classifiers consume the confirmatory markers whose
        unavailability is the actual bottleneck to early detection. There is
        little work establishing how much risk discrimination survives when
        those markers are withheld and only routine checkup measurements
        remain.
  \item \textbf{Symmetric treatment of asymmetric errors.} Reported CKD models
        are optimised and reported on aggregate accuracy, without explicit
        treatment of the high-to-low misclassification as the dominant clinical
        harm, and without an architecture that provides a natural place to
        control it.
  \item \textbf{Risk scores without communication.} Screening models terminate
        at a numeric output. Personalised, guideline-grounded, readable
        communication of what that output means---and its evaluation---is
        largely absent.
  \item \textbf{Unmeasured, unattributed LLM safety.} Where language models are
        used for health education, safety is usually asserted through prompt
        design rather than measured, and where it is measured, the measurement
        does not isolate \emph{which} component produces it. Consequently there
        is no evidence base for deciding whether fine-tuning, retrieval, or an
        explicit guardrail is the mechanism worth investing in.
\end{enumerate}

This thesis addresses gap~(1) by constructing a routine-input-only feature set
with a KDIGO-derived label; gap~(2) through a two-stage cascade with
independently tuned thresholds and explicit reporting of the high-to-low error
count; gap~(3) by building and deploying a grounded education layer; and
gap~(4) through a pre-registered six-arm ablation that measures and attributes
safety behaviour component by component.
